
\documentclass{article}
\usepackage{colortbl}
\usepackage{makecell}
\usepackage{multirow}
\usepackage{supertabular}

\begin{document}

\newcounter{utterance}

\centering \large Interaction Transcript for game `referencegame', experiment `number\_grids', episode 13 with reflexion{-}Qwen3.5{-}9B.
\vspace{24pt}

{ \footnotesize  \setcounter{utterance}{1}
\setlength{\tabcolsep}{0pt}
\begin{supertabular}{c@{$\;$}|p{.15\linewidth}@{}p{.15\linewidth}p{.15\linewidth}p{.15\linewidth}p{.15\linewidth}p{.15\linewidth}}
    \# & $\;$A & \multicolumn{4}{c}{Game Master} & $\;\:$B\\
    \hline

    \theutterance \stepcounter{utterance}  
    & & \multicolumn{4}{p{0.6\linewidth}}{
        \cellcolor[rgb]{0.9,0.9,0.9}{
            \makecell[{{p{\linewidth}}}]{
                \texttt{\tiny{[P1$\langle$GM]}}
                \texttt{You have attempted this same task before and did not succeed. These are the plans you wrote after those attempts.} \\
\texttt{Trial \#1: New plan: My previous description was too vague and failed to explicitly distinguish the target from the distractors by pinpointing the exact row differences. Specifically, I mentioned a "complete rectangular border" which is true for the target, but I did not clearly state that the distractors have a "broken" bottom border (missing the bottom{-}left cell) or a "broken" middle section in the third row. Next time, I will construct the expression by first confirming the shared features (the vertical line in column 2 and the top/bottom rows in the target) and then immediately contrasting them with the specific missing cells in the distractors: explicitly stating that the target has filled cells at (row 3, col 1) and (row 4, col 1) which are empty in both distractors, and that the target has a solid bottom row while Distractor 1 has a gap at the bottom{-}left and Distractor 2 has a gap at the bottom{-}left and a different middle pattern. I will ensure the phrase "unlike the distractors" is followed by the precise coordinates or relative positions of the differing cells.} \\
\\ 
\texttt{Trial \#2: New plan: My previous description failed because it relied on vague terms like "solid bottom{-}left corner" and "broken border" without explicitly mapping the specific cell coordinates that differ between the grids. The target grid is unique because it has filled cells at (row 3, col 1) and (row 4, col 1), whereas Distractor 1 has a filled cell at (row 3, col 3) and Distractor 2 has an entirely empty fourth row. Next time, I will construct the expression by first stating the shared vertical line in column 2, then immediately contrasting the target's specific filled cells at (3,1) and (4,1) against the specific differing features of the distractors: explicitly noting that Distractor 1 has a filled cell at (3,3) instead of (3,1), and Distractor 2 has an empty row 4 instead of the filled cells at (4,1). I will ensure every distinguishing feature is tied to exact row and column coordinates to eliminate ambiguity.} \\
\\ 
\texttt{Trial \#3: New plan: My previous attempts failed because I described the target's unique features (the left border) but incorrectly attributed the distractors' differences to specific coordinates that were either wrong or insufficient to guarantee uniqueness. Specifically, I claimed Distractor 1 had a filled cell at (3,3), but looking closely, Distractor 1 actually has a filled cell at (3,3) while the target has an empty cell there; however, my description of Distractor 2 as having an "entirely empty fourth row" was factually incorrect (Distractor 2 has filled cells in row 4, specifically at cols 2,3,4,5, but the key difference is the missing cell at (4,1)). The core mistake was failing to explicitly state that the target is the *only* grid with filled cells at both (3,1) and (4,1). Next time, I will construct the expression by first defining the common structure (vertical line in col 2, full top/bottom rows), then immediately asserting that the target is uniquely identified by having filled cells at (row 3, col 1) and (row 4, col 1), while explicitly stating that both Distractor 1 and Distractor 2 have empty cells at these exact coordinates, thereby creating a definitive logical exclusion.} \\
\\ 
\\ 
\texttt{You are given three grids, where each of them is 5 by 5 in size.} \\
\texttt{Grids have empty cells marked with "▢" and filled cells marked with "X".} \\
\texttt{Your task is to generate a referring expression that best describes the target grid while distinguishing it from the two other distractor grids.} \\
\texttt{The first grid is the target grid, and the following two grids are the distractors.} \\
\\ 
\texttt{Target grid:} \\
\\ 
\texttt{□ X X X X} \\
\texttt{□ X □ □ X} \\
\texttt{□ X □ □ X} \\
\texttt{□ X □ □ X} \\
\texttt{□ X X X X} \\
\\ 
\texttt{Distractor grid 1:} \\
\\ 
\texttt{□ X X X X} \\
\texttt{□ X □ □ X} \\
\texttt{□ X X X X} \\
\texttt{□ X □ □ X} \\
\texttt{□ X X X X} \\
\\ 
\texttt{Distractor grid 2:} \\
\\ 
\texttt{□ X X X X} \\
\texttt{□ X □ □ X} \\
\texttt{□ X X X X} \\
\texttt{□ □ □ □ □} \\
\texttt{□ X X X X} \\
\\ 
\texttt{Instruction: Describe the target grid.} \\
\texttt{Generate the referring expression starting with the tag "Expression: " for the given target grid. Omit any other text.} \\
            }
        }
    }
    & \\ \\

    \theutterance \stepcounter{utterance}  
    & \multicolumn{4}{p{0.6\linewidth}}{
        \cellcolor[rgb]{0.9,0.9,0.9}{
            \makecell[{{p{\linewidth}}}]{
                \texttt{\tiny{[P1$\rangle$GM]}}
                \texttt{Expression: The grid containing a vertical line of filled cells in the second column, with the top and bottom rows completely filled, and uniquely distinguished by having filled cells at both (row 3, column 1) and (row 4, column 1), whereas the other two grids have empty cells at these specific coordinates.} \\
            }
        }
    }
    & & \\ \\

    \theutterance \stepcounter{utterance}  
    & & & \multicolumn{2}{p{0.3\linewidth}}{
        \cellcolor[rgb]{0.9,0.9,0.9}{
            \makecell[{{p{\linewidth}}}]{
                \texttt{\tiny{[GM$|$GM]}}
                \texttt{Expression: The grid containing a vertical line of filled cells in the second column, with the top and bottom rows completely filled, and uniquely distinguished by having filled cells at both (row 3, column 1) and (row 4, column 1), whereas the other two grids have empty cells at these specific coordinates.} \\
            }
        }
    }
    & & \\ \\

    \theutterance \stepcounter{utterance}  
    & & \multicolumn{4}{p{0.6\linewidth}}{
        \cellcolor[rgb]{0.9,0.9,0.9}{
            \makecell[{{p{\linewidth}}}]{
                \texttt{\tiny{[GM$\rangle$P2]}}
                \texttt{You have attempted this same task before and did not succeed. These are the plans you wrote after those attempts.} \\
\texttt{Trial \#1: New plan: I incorrectly identified the "First" grid because I focused on the vertical line in the second column but failed to verify the horizontal border requirement. The expression demands a complete rectangular border (filled first and last rows) with an empty interior except for the second column. The "First" grid has empty cells in the first and last rows, violating the border condition. The "Second" grid has a filled first row but an empty last row. The "Third" grid has a filled first row, a filled last row, and a vertical line in the second column, while the interior (rows 2{-}4, columns 1, 3{-}5) is empty. Next time, I will strictly check for the presence of filled cells in both the first and last rows before confirming the vertical line, ensuring the complete rectangular border is present. The correct selection is the Third grid.} \\
\texttt{Answer: Third} \\
\\ 
\texttt{Trial \#2: New plan: I failed because I misinterpreted the expression's description of the "distractors" versus the "target." The expression explicitly states that the target has filled cells at (row 3, col 1) and (row 4, col 1), which are empty in the distractors. Upon re{-}evaluating the grids:} \\
\texttt{{-} **First Grid**: Has filled cells at (3,1) and (4,1). It also has a complete vertical line in column 2 and filled top/bottom rows. This matches the target description perfectly.} \\
\texttt{{-} **Second Grid**: Has empty cells at (3,1) and (4,1). This matches the description of a distractor.} \\
\texttt{{-} **Third Grid**: Has empty cells at (3,1) and (4,1). This also matches the description of a distractor.} \\
\texttt{My previous error was selecting the Third grid, likely due to a hasty visual scan of the bottom rows without verifying the specific coordinates mentioned in the text. Next time, I will strictly cross{-}reference the specific coordinate claims in the expression (specifically rows 3 and 4, column 1) against the grid layouts before making a selection. The correct selection is the First grid.} \\
\texttt{Answer: First} \\
\\ 
\texttt{Trial \#3: New plan: I failed because I misread the coordinate descriptions in the expression. The text explicitly states the target has filled cells at (row 3, col 1) and (row 4, col 1), which corresponds to the "First" grid. However, I incorrectly assumed the "Third" grid matched the description of having a complete rectangular border, ignoring that the expression defines the target's unique feature as the filled left border cells in rows 3 and 4, while the distractors (Second and Third) are defined by having empty cells at those specific coordinates. Next time, I will strictly map the coordinate claims (row 3, col 1 and row 4, col 1) to the visual grids first to identify the target, rather than relying on a general visual scan of borders. The correct selection is the First grid.} \\
\texttt{Answer: First} \\
\\ 
\\ 
\texttt{You are given three grids, where each of them is 5 by 5 in size.} \\
\texttt{Grids have empty cells marked with "▢" and filled cells marked with "X".} \\
\texttt{You are also given a referring expression that describes one of the given grids.} \\
\texttt{Your task is to select the grid that matches the given referring expression.} \\
\texttt{Generate only the number (in text) of the grid that matches the given expression by selecting first, second, or third.} \\
\\ 
\texttt{First:} \\
\\ 
\texttt{□ X X X X} \\
\texttt{□ X □ □ X} \\
\texttt{□ X X X X} \\
\texttt{□ X □ □ X} \\
\texttt{□ X X X X} \\
\\ 
\texttt{Second:} \\
\\ 
\texttt{□ X X X X} \\
\texttt{□ X □ □ X} \\
\texttt{□ X □ □ X} \\
\texttt{□ X □ □ X} \\
\texttt{□ X X X X} \\
\\ 
\texttt{Third:} \\
\\ 
\texttt{□ X X X X} \\
\texttt{□ X □ □ X} \\
\texttt{□ X X X X} \\
\texttt{□ □ □ □ □} \\
\texttt{□ X X X X} \\
\\ 
\texttt{Expression: The grid containing a vertical line of filled cells in the second column, with the top and bottom rows completely filled, and uniquely distinguished by having filled cells at both (row 3, column 1) and (row 4, column 1), whereas the other two grids have empty cells at these specific coordinates.} \\
\texttt{Question: Which grid does the expression refer to?} \\
\texttt{Start with the tag "Answer: ", followed by your selection. Omit any other text.} \\
            }
        }
    }
    & \\ \\

    \theutterance \stepcounter{utterance}  
    & & & \multicolumn{4}{p{0.6\linewidth}}{
        \cellcolor[rgb]{0.9,0.9,0.9}{
            \makecell[{{p{\linewidth}}}]{
                \texttt{\tiny{[GM$\langle$P2]}}
                \texttt{Answer: First} \\
            }
        }
    }
     \\ \\

    \theutterance \stepcounter{utterance}  
    & & & \multicolumn{2}{p{0.3\linewidth}}{
        \cellcolor[rgb]{0.9,0.9,0.9}{
            \makecell[{{p{\linewidth}}}]{
                \texttt{\tiny{[GM$|$GM]}}
                \texttt{first} \\
            }
        }
    }
    & & \\ \\

    \theutterance \stepcounter{utterance}  
    & \multicolumn{4}{p{0.6\linewidth}}{
        \cellcolor[rgb]{0.95,0.95,0.95}{
            \makecell[{{p{\linewidth}}}]{
                \texttt{\tiny{[P2$\rightarrow$P1]}}
                \texttt{Post{-}episode Reflexion (not sent to GM):} \\
\texttt{New plan: My previous attempts failed because I focused on describing the target's unique features in isolation rather than explicitly contrasting them against the specific, incorrect features present in the distractors. Specifically, I incorrectly claimed Distractor 1 had a filled cell at (3,3) while the target did not, but the target actually has an empty cell at (3,3) and the distractor has a filled one; however, the definitive differentiator is that the target is the *only* grid with filled cells at both (3,1) and (4,1). Next time, I will construct the expression by first stating the common structure (vertical line in column 2, full top/bottom rows), then immediately asserting that the target is uniquely identified by having filled cells at (row 3, col 1) and (row 4, col 1), while explicitly stating that both Distractor 1 and Distractor 2 have empty cells at these exact coordinates, thereby creating a definitive logical exclusion without making false claims about other cells.} \\
            }
        }
    }
    & & \\ \\

    \theutterance \stepcounter{utterance}  
    & & & \multicolumn{4}{p{0.6\linewidth}}{
        \cellcolor[rgb]{0.95,0.95,0.95}{
            \makecell[{{p{\linewidth}}}]{
                \texttt{\tiny{[P1$\rightarrow$P2]}}
                \texttt{Post{-}episode Reflexion (not sent to GM):} \\
\texttt{New plan: I failed because I misinterpreted the coordinate system, likely confusing 0{-}indexed and 1{-}indexed rows or miscounting the visual lines. The expression explicitly requires filled cells at (row 3, column 1) and (row 4, column 1). In the provided grids, the "First" grid has filled cells at these exact positions (the third and fourth rows down in the first column), while the "Second" and "Third" grids have empty cells there. My previous error was over{-}prioritizing the "complete rectangular border" description, which is actually a distractor feature in this specific puzzle setup, rather than focusing on the unique coordinate constraint that definitively identifies the target. Next time, I will strictly map the specific coordinate claims (row 3, col 1 and row 4, col 1) to the visual grids first to identify the target, ignoring general visual patterns like borders that do not match the unique coordinate constraint. The correct selection is the First grid.} \\
\texttt{Answer: First} \\
            }
        }
    }
     \\ \\

\end{supertabular}
}

\end{document}
